% Minimal monthly report - Pressa public template (MIT). https://pressa.dev/templates/monthly-report?design=minimal-monthly-report
\documentclass[10pt,a4paper]{article}
\usepackage[T1]{fontenc}
\usepackage[scaled]{helvet}
\renewcommand{\familydefault}{\sfdefault}
\usepackage[margin=18mm,top=16mm,bottom=20mm]{geometry}
\usepackage[table]{xcolor}
\usepackage{tabularx}
\usepackage{array}
\usepackage{longtable}
\usepackage{booktabs}
\usepackage{enumitem}
\usepackage{titlesec}
\usepackage{pgfplots}
\usepackage{fancyhdr}
\usepackage[hidelinks]{hyperref}
\pgfplotsset{compat=1.18}

% `accent` is the colour a visitor can change; tints mix it inline (accent!8),
% never through \colorlet, so they follow the change.
\definecolor{accent}{HTML}{0F766E}
\definecolor{ink}{HTML}{1F2937}
\definecolor{muted}{HTML}{6B7280}
\color{ink}
\setlength{\parindent}{0pt}
\setlength{\parskip}{5pt}
\titleformat{\section}{\normalsize\bfseries}{}{0pt}{}
\titlespacing*{\section}{0pt}{18pt}{6pt}
\newcommand{\fieldname}[1]{\textcolor{muted}{\scriptsize\MakeUppercase{#1}}}
\setlist[itemize]{leftmargin=1.2em, itemsep=2pt, topsep=2pt, label={\color{accent}\textbullet}}
\newcolumntype{C}{>{\centering\arraybackslash}X}
\pagestyle{fancy}
\fancyhf{}
\renewcommand{\headrulewidth}{0pt}
\fancyfoot[L]{\small\color{muted} {{ company }} \quad {{ report_title }}, {{ period }}}
\fancyfoot[R]{\small\color{muted}\thepage}

\begin{document}
\fieldname{ {{- company -}} }\par\vspace{4pt}
{\fontsize{26}{30}\selectfont\mbox{}{{ report_title }}}\par\vspace{6pt}
{\small\color{muted}\mbox{}{{ period }} \quad\textbullet\quad Prepared by {{ prepared_by }}}\par\vspace{8pt}
\textcolor{accent}{\rule{12mm}{2pt}}

\vspace{10pt}
{{ summary }}
{% if kpis.size > 0 %}
\vspace{8pt}
{\renewcommand{\arraystretch}{1.25}\setlength{\tabcolsep}{9pt}
% One column per KPI, so a fifth or a third KPI still builds.
\noindent\begin{tabularx}{\linewidth}{@{}{% for kpi in kpis %}C{% unless forloop.last %}!{\color{black!12}\vrule width 0.5pt}{% endunless %}{% endfor %}@{}}
{% for kpi in kpis %}\vspace{4pt}{\scriptsize\color{muted}\MakeUppercase{ {{- kpi.label -}} }\par}\vspace{3pt}{\fontsize{20}{23}\selectfont {{ kpi.value }}\par}\vspace{1pt}{\small\color{accent} {{ kpi.change }}\par}\vspace{4pt}{% unless forloop.last %} & {% endunless %}{% endfor %} \\
\end{tabularx}}
{% endif %}

{% if chart.size > 0 %}\section{ {{- chart_title -}} }
\begin{tikzpicture}
\begin{axis}[
  ybar, width=\linewidth, height=5.6cm, bar width=15pt, ymin=0,
  axis lines*=left, axis line style={draw=muted!40}, ymajorgrids, grid style={draw=muted!15},
  enlarge x limits=0.06, xtick=data, tick style={draw=none},
  tick label style={font=\small, color=muted},
  scaled y ticks=false, /pgf/number format/assume math mode=true,
  yticklabel style={/pgf/number format/fixed},
  nodes near coords, nodes near coords style={font=\scriptsize, color=ink},
  symbolic x coords={ {% for point in chart %}{{ point.label }}{% unless forloop.last %},{% endunless %}{% endfor %} },
]
\addplot[fill=accent, draw=none] coordinates { {% for point in chart %}({{ point.label }},{{ point.value | round: 2 }}) {% endfor %} };
\end{axis}
\end{tikzpicture}
{% endif %}

{% if table.size > 0 %}\section{ {{- table_title -}} }
{\renewcommand{\arraystretch}{1.6}\arrayrulecolor{black!12}
% A longtable with fixed columns: a long table breaks across pages with its
% header repeated, in one compile pass.
\newlength{\metriccol}\setlength{\metriccol}{24mm}\setlength{\LTpre}{0pt}\setlength{\LTpost}{0pt}\setlength{\LTleft}{0pt}\setlength{\LTright}{0pt}
\begin{longtable}{@{}>{\raggedright\arraybackslash}p{\dimexpr\linewidth-6\tabcolsep-3\metriccol\relax} >{\raggedleft\arraybackslash}p{\metriccol} >{\raggedleft\arraybackslash}p{\metriccol} >{\raggedleft\arraybackslash}p{\metriccol}@{}}
\fieldname{Metric} & \fieldname{This period} & \fieldname{Previous} & \fieldname{Change} \\
\hline
\endhead
{% for row in table %}{{ row.metric }} & {{ row.current }} & {{ row.previous }} & {{ row.change }} \\
\hline
{% endfor %}\end{longtable}}
{% endif %}

\begin{minipage}[t]{0.48\linewidth}
{% if highlights.size > 0 %}\section{Highlights}
\begin{itemize}
{% for item in highlights %}  \item {{ item }}
{% endfor %}\end{itemize}{% endif %}
\end{minipage}\hfill
\begin{minipage}[t]{0.48\linewidth}
{% if next_steps.size > 0 %}\section{Next steps}
\begin{itemize}
{% for item in next_steps %}  \item {{ item }}
{% endfor %}\end{itemize}{% endif %}
\end{minipage}

\end{document}
